\clearpage
\section{Dereference coherence and the complete store interface}
\label{sec:copying-dereference-coherence}
Every actual contextual read now has a unique complete successor up to
recursive state isomorphism. Together with the preceding allocation,
assignment and equality results, this proves successor coherence for all
four store-operation heads, including both equality branches. The proof
uses the original Figure~19(b) rule, described in
\href{https://arxiv.org/html/1907.01257v5\#S10.SS1}{HTML Section~10.1}.
It retains the existing store and atom tree, consumes the instance by weakening,
and connects the store and request to the exact distributor
$D^{\mathrm{term}}_{1,2}$, including its weakening and two contractions.

The key quantitative conclusion is a continuation law. With arbitrary contents
attached, a read succeeds in $n$ transitions if and only if $n=5+r$ and
the resulting copying request succeeds in $r$ transitions. In particular,
the five-step read prefix is not a completed evaluation. This strengthens
the earlier chosen local trace to a statement about every successful execution
of the full combined machine, including all contextual matching choices.

The dependency chain is: unique source read match; a finite read graft;
typed target exchange; commutation with contexts; explicit isomorphism to
the original reduct; induction over actual rewrites; successor coherence;
exact cost decomposition. The graft is solely a proof construction. It does
not add a transition or change the terms of Proposition~12.2, the translation
of Figure~16, or any stability condition in the operational rules.

This completes the store-operation part of Appendix~B.1's successor obligation.
Section~\ref{sec:copying-arithmetic-coherence} subsequently completes arithmetic.
Section~\ref{sec:copying-application-coherence} completes application. Copying still requires successor comparison, and
the complete machine's refocusing and specimen laws remain open. The four
broad packages stay open. The mathematical-effort estimate remains roughly
50--65\%, with low confidence: these proofs do not complete the template
argument or the final observational equivalence.

\Needspace{12\baselineskip}
\subsection{A finite graft for the exact distributor}

\termdef{read-graft}
The \emph{read graft} is a graph with one term input $q$, one term output $i$,
three fresh internal term ports $a,b,c$ and four edges. Its first two edges
are weakenings and its last two are contractions. The ordered source and
target lists are
\[
\begin{array}{c|c|c}
\text{edge}&\text{sources}&\text{targets}\\\hline
0&[]&[i]\\
1&[]&[b]\\
2&[b,c]&[a]\\
3&[a,q]&[c].
\end{array}
\]
Edge zero will discard the consumed instance. Edges one to three are the
literal $D^{\mathrm{term}}_{1,2}$ pattern, temporarily joining its first input
to its output at $a$. That temporary feedback is removed by target exchange
with the store-content port before the graph is compared with a read result.
No claim of copyability or stability is made for the intermediate graft.

\lean{Copying/ReadStub}{12}{13}{ReadStub.label}

\lean{Copying/ReadStub}{15}{17}{ReadStub.source}

\lean{Copying/ReadStub}{19}{21}{ReadStub.target}

Every index of the four-element edge type is one of $0,1,2,3$, by its
finite bound. Enumerating these cases verifies absence of repeated ports
in each list and uniqueness of each source and target incidence. For coverage,
$q,a,b,c$ each occur once as a source and $i,a,b,c$ each occur once as a
target; $i$ is not a source and $q$ is not a target. All labels are term,
so the listed arities also verify typing. These facts construct the finite
primitive graph. Relabelling its primitive edges at any depth $d$ gives
the corresponding hypernet, with the same carriers and incidences.

\lean{Copying/ReadStub}{9}{9}{ReadStub.fin_four}

\lean{Copying/ReadStub}{24}{75}{ReadStub.base}

\lean{Copying/ReadStub}{77}{78}{ReadStub.net}

\Needspace{12\baselineskip}
\subsection{State-level transport of target exchange}

A non-input port remains a non-input under any incidence-preserving match
embedding: target coverage supplies an edge targeting it, the embedding
transports that incidence, and no target in the receiving graph can be an
external input. This argument applies even when boundary attachment maps
are not globally injective.

\lean{Copying/TargetSwapStates}{9}{12}{TargetSwap.embedded_not_input}

The state-level target exchange $S_{a,b}(G,v)$ consists of the exchanged
graph and the search token at the unchanged port $v$. Only target incidences
are exchanged, so moving the token by the exchange would be incorrect.
If two graphs are equal, transporting their ports along that equality
preserves both the non-input conditions and equality of endpoint labels.
Eliminating the equality reduces each assertion to its original proof.
The same elimination gives the state isomorphism back to the original
exchanged state.

\lean{Copying/TargetSwapStates}{14}{16}{TargetSwap.state}

\lean{Copying/TargetSwapStates}{19}{21}{TargetSwap.cast_not_input}

\lean{Copying/TargetSwapStates}{23}{27}{TargetSwap.cast_typed}

\lean{Copying/TargetSwapStates}{29}{35}{TargetSwap.state_cast}

Suppose a recursive graph isomorphism maps $a,b,v$ to $a',b',w$.
The earlier target-exchange renaming theorem supplies an isomorphism of
the exchanged graphs, and the given token equation supplies the state
isomorphism. For plugging, use the earlier exchange-plugging graph isomorphism.
Its carrier map is the identity and its insertion-map equation identifies
the token. Thus exchanging the embedded endpoints in the whole graph agrees
with inserting the exchanged component state.

\lean{Copying/TargetSwapStates}{38}{45}{TargetSwap.state_congr}

\lean{Copying/TargetSwapStates}{48}{61}{TargetSwap.state_plug}

For simultaneous expansion, let $F$ and $K$ agree except at the selected
component. Assume that exchanging the selected endpoints in $F(e)$ gives
a state isomorphic to the specified state in $K(e)$. First commute exchange
through the expansion, obtaining the family $P$ which exchanges only $F(e)$.
At $e$, the family definition identifies $P(e)$ with that exchanged graph;
elsewhere $P$ agrees with $F$, hence with $K$. Apply expansion congruence to
the given component state isomorphism. The insertion-map equation tracks the
token through both comparisons. Other components may contain arbitrary
recursive boxes and cycles.

\lean{Copying/TargetSwapStates}{64}{91}{TargetSwap.state_expand_result}

\Needspace{12\baselineskip}
\subsection{Unique source data and the canonical read successor}

\termdef{dereference-match}
A \emph{dereference match} at a port $v$ records a dereference head $e$, a
store $s$, its content port $o$, and an atom-trace length $k$. Its evidence
says that $v$ is a source of $e$, $s$ has the store label and singleton target
$[o]$, and operand zero of the dereference resolves to $s$ in $k$ contractions.
This is a description of source incidences and resolution, not an assumed
result or matching oracle.

\lean{Copying/DereferenceMatch}{8}{16}{Dereference.Match}

Match embeddings transport the head, store and content port, preserving
the trace length. Ordered target transport preserves $[o]$, and operand
transport preserves store resolution. Matches at the same source request
are unique: source-head uniqueness identifies the heads; uniqueness of
operand resolution identifies the stores and trace lengths; the singleton
target equations identify the content ports. Proof irrelevance identifies
the remaining evidence. Because $o$ is a store target, it is not an input;
the store's target-typing law gives it type term.

\lean{Copying/DereferenceMatch}{20}{29}{Dereference.Match.map}

\lean{Copying/DereferenceMatch}{31}{35}{Dereference.Match.unique}

\lean{Copying/DereferenceMatch}{37}{39}{Dereference.Match.contents_not_input}

\lean{Copying/DereferenceMatch}{41}{45}{Dereference.Match.contents_typed}

The match's filler is the read graft transported to the head's exact
interface using its primitive label. Replacing just $e$ by this filler
defines the grafted graph. Its fresh port is the graft's internal port zero,
and its retained port is the host image of $o$. The fresh port is internal
by construction. The retained port is non-input by the original property
and the three cases of the host port map. Both ports have type term: this
is immediate for the graft port, and follows from preservation of host
labels for the content port. Thus target exchange is well typed.

\lean{Copying/DereferenceMatch}{47}{51}{Dereference.Match.filler}

\lean{Copying/DereferenceMatch}{53}{54}{Dereference.Match.grafted}

\lean{Copying/DereferenceMatch}{56}{56}{Dereference.Match.fresh}

\lean{Copying/DereferenceMatch}{57}{58}{Dereference.Match.retained}

\lean{Copying/DereferenceMatch}{60}{61}{Dereference.Match.fresh_not_input}

\lean{Copying/DereferenceMatch}{63}{65}{Dereference.Match.retained_not_input}

\lean{Copying/DereferenceMatch}{67}{71}{Dereference.Match.endpoints_typed}

\termdef{canonical-read-successor}
The \emph{canonical read successor} is the search state obtained by exchanging
the targets at those fresh and retained ports in the grafted graph, retaining
the old request token. The store, previously targeting $o$, now targets $a$;
the distributor root, previously targeting $a$, now targets $o$. Its ordered
inputs therefore connect the store and resumed request to the original stored
contents. The separate weakening connects to the consumed instance.

\lean{Copying/DereferenceMatch}{74}{76}{Dereference.Match.canonical}

\Needspace{11\baselineskip}
Every original local read source supplies a match. Its head and store are
their explicit expanded occurrences; the content port is the source graph's
single term output. The earlier finite store-resolution theorem supplies a
trace length and operand proof. Classical choice selects this witness, while
match uniqueness proves that its choice cannot change the canonical successor.

\lean{Copying/DereferenceMatch}{79}{88}{Dereference.localMatch}

\Needspace{12\baselineskip}
\subsection{Commutation with renaming and contexts}

A recursive isomorphism of source graphs induces an isomorphism of their
grafted graphs. Use the carrier bijections supplied by host plugging. On
retained edges, labels and incidence lists obey the original isomorphism.
On the four inserted edges, labels are identical; the insertion-map equation,
applied portwise to each ordered list, proves both incidence equations.
Boundary casts are checked separately on internal, input and output ports.
The fresh port maps to the fresh port, and host-port naturality transports
both the content port and the token. State-level exchange congruence now
proves renaming compatibility of the canonical successor.

\lean{Copying/DereferenceCanonicalTransport}{8}{52}{Dereference.Match.graft_rename}

\lean{Copying/DereferenceCanonicalTransport}{54}{61}{Dereference.Match.canonical_rename}

For plugging, associativity of replacement first identifies grafting at
the embedded head with inserting the grafted component. Its middle-component
port equation identifies the retained content port and request token, and
the fresh graft port maps directly. The two embedded endpoints remain
non-inputs, and insertion preserves their labels. State-level exchange
congruence followed by exchange-plugging gives the desired canonical-state
comparison for any typed host.

\lean{Copying/DereferenceCanonicalTransport}{62}{81}{Dereference.Match.canonical_plug}

For simultaneous expansion, replace the head only in the selected component
of a family $F$, giving a family $P$. Its selected graph equals the grafted
component, so equality transport supplies its endpoint conditions and token.
The expansion-plugging isomorphism compares whole-graph grafting with
expansion by $P$; its inserted-port and component-port equations identify
the fresh endpoint, retained endpoint and token. Apply exchange congruence,
then the state-level exchange-expansion theorem with the given component
successor comparison. This allows the new family $K$ to use an independently
represented successor, as the actual rewrite relation requires.

\lean{Copying/DereferenceCanonicalTransport}{83}{118}{Dereference.Match.canonical_expand_result}

\Needspace{12\baselineskip}
\subsection{The explicit isomorphism to Figure 19(b)}

The abbreviation \texttt{replaced} names the graph of the canonical
successor for the original local source. Its forward internal-vertex map
retains the three old frame vertices and every atom-tree vertex. The graft's
fresh vertex becomes the fourth frame vertex of the original reduct, and
its other two vertices become the distributor's two internal vertices.
\lean{Copying/DereferenceReplacement}{8}{9}{Dereference.Replacement.replaced}

\lean{Copying/DereferenceReplacement}{11}{18}{Dereference.Replacement.vertexTo}

\Needspace{10\baselineskip}
The inverse separates these cases and restores the replacement tags.
Primitive components have no internal vertices. Enumerating the finite
frame and graft indices proves that the two maps are inverse bijections.

\lean{Copying/DereferenceReplacement}{20}{27}{Dereference.Replacement.vertexFrom}

\lean{Copying/DereferenceReplacement}{29}{41}{Dereference.Replacement.vertices}

\Needspace{12\baselineskip}
The edge map retains the instance, store and every atom-tree occurrence.
The deleted head is excluded by the retained-edge subtype. Graft edge zero
becomes the discarded-instance weakening; graft edges one, two and three
become respectively the distributor's weakening and its two contractions.
The inverse restores the finite indices and retained-edge tags. Enumerating
the four graft indices and the one-edge primitive components proves both
inverse equations and gives the edge bijection.

\lean{Copying/DereferenceReplacement}{43}{49}{Dereference.Replacement.edgeTo}

\lean{Copying/DereferenceReplacement}{51}{57}{Dereference.Replacement.edgeFrom}

\lean{Copying/DereferenceReplacement}{59}{71}{Dereference.Replacement.edges}

The vertex bijection preserves every atom-tree port's attachment. Internal
tree ports and the single output follow directly from the definitions; for
inputs, split the first input from the remaining atom-interface inputs and
reduce the respective frame attachments. Moreover, target exchange fixes
every such tree port: retained host ports cannot be the fresh graft port,
and none of the tree's attached internal, input or output ports is the
external content port. These observations also cover noncanonical atom trees
with any number of additional references at their boundary.

\lean{Copying/DereferenceReplacement}{74}{86}{Dereference.Replacement.map_tree}

\lean{Copying/DereferenceReplacement}{90}{111}{Dereference.Replacement.exchange_tree}

These bijections preserve all labels and ordered incidence lists, giving
an isomorphism of complete graphs. Tree sources use the attachment equation;
tree targets additionally use the fact that exchange fixes them. The instance
target is unchanged. The store target moves from the old content port to
the fresh port. In the graft, the distributor root target makes the opposite
move, whereas the discarded-instance weakening, distributor weakening and
inner contraction targets are fixed. All remaining source and label cases
are the displayed finite incidences. The request input is fixed, so this
graph isomorphism lifts to a state isomorphism with the original resumed
search state. This establishes correspondence with the unchanged reduct,
not merely agreement at the store endpoint.

\lean{Copying/DereferenceReplacement}{114}{186}{Dereference.Replacement.iso}

\lean{Copying/DereferenceReplacement}{188}{190}{Dereference.Replacement.canonical_local}

\Needspace{12\baselineskip}
\subsection{Coherence of every actual contextual read}

For a local step admitting a dereference match, inspect the rule constructor.
Every non-read constructor has a different active head label, contradicting
source-head uniqueness. In the read case, uniqueness identifies the given
match with the original local match, and the explicit isomorphism applies.
The parallel local inversion theorem extracts a match from any actual local
step at a dereference head, using the same constructor analysis.

\lean{Copying/DereferenceCoherence}{9}{31}{Combined.local_dereference_canonical}

\lean{Copying/DereferenceCoherence}{35}{54}{Combined.local_dereference_match}

Inducting over actual rewrite derivations extends match extraction to all
contextual reads. At a local step use inversion. Under recursive renaming,
plugging or simultaneous expansion, reflect the active head into the rewritten
component, extract its match by induction, and transport it outward through
the corresponding match embedding. No assumption of unique decomposition
or a pre-existing deterministic matcher is used.

\lean{Copying/DereferenceCoherence}{58}{74}{Combined.rewrite_dereference_match}

A second induction compares each actual read successor with the canonical
successor of its source match. At a contextual constructor, extract the inner
match and identify its outward transport with the supplied match by uniqueness.
Apply induction inside the component and then the appropriate canonical
renaming, plugging or expansion theorem. In the renaming case, compose with
the separately specified target isomorphism. Thus independent contextual
decompositions are compared through one source-determined state.

\lean{Copying/DereferenceCoherence}{77}{102}{Combined.rewrite_dereference_canonical}

Given two reads at the same request, extract a match from the first and
compare both targets with its canonical successor. Composing one comparison's
inverse with the other gives a complete recursive state isomorphism between
the targets. For full-machine steps in rewrite mode, intrinsic search is
impossible, so both steps must be rewrites and the same theorem applies.
The generic coherent-step cost law then shows that one actual read step
consumes exactly one transition and leaves the same full spectrum of future
successful costs, independently of matching choices.

\lean{Copying/DereferenceCoherence}{105}{112}{Combined.dereference_rewrite_deterministic}

\lean{Copying/DereferenceCoherence}{115}{124}{Combined.dereference_step_deterministic}

\lean{Copying/DereferenceCoherence}{127}{131}{Combined.dereference_cost_spectrum}

An unbounded step is a lift of a bounded step. Destructing two such lifts
and applying bounded coherence gives the proved representation relation
between their successors. Representation invariance gives observational
equivalence at every cost preorder, including equality of costs. These
statements impose no finite-depth cutoff and no global determinism premise.

\lean{Copying/DereferenceCoherence}{134}{143}{Combined.dereference_unbounded_deterministic}

\lean{Copying/DereferenceCoherence}{146}{149}{Combined.dereference_unbounded_equiv}

\Needspace{12\baselineskip}
\subsection{All store operations under one theorem}

\termdef{store-successor-coherence}
The predicate \texttt{StoreHead} holds when the selected request is an
allocation, dereference, assignment or atom-equality head. Its four constructors
retain the actual head-incidence evidence. Equality includes both the same-cell
and distinct-cell operational rules. It asserts the kind of head, not that
every such source admits a rewrite.

\lean{Copying/StoreCoherence}{8}{12}{Combined.StoreHead}

\emph{Store successor coherence} says that any two actual rewrites from
one of these heads have recursively isomorphic complete successors, including
the resumed token. Case analysis on \texttt{StoreHead} applies the four
independently proved family theorems. The full-machine step theorem follows
by the same case analysis. This closes the store-operation portion of
successor coherence while keeping existence and stable-payload conditions
in the original operational rules.

\lean{Copying/StoreCoherence}{15}{22}{Combined.store_rewrite_deterministic}

\lean{Copying/StoreCoherence}{25}{32}{Combined.store_step_deterministic}

The generic coherent-step theorem now yields a uniform cost law:
if $s\to t$ is any actual store computation, then $s$ terminates in $n$
exactly when $n=1+r$ and $t$ terminates in $r$. Finally, destructing unbounded
lifts transfers coherence to represented states, and representation invariance
gives successor equivalence for every cost preorder. These are genuine
consequences of the proved store comparisons; they do not supply the still
then-missing application, arithmetic or copying comparisons.

\lean{Copying/StoreCoherence}{35}{38}{Combined.store_cost_spectrum}

\lean{Copying/StoreCoherence}{41}{49}{Combined.store_unbounded_deterministic}

\lean{Copying/StoreCoherence}{52}{55}{Combined.store_unbounded_equiv}

\Needspace{12\baselineskip}
\subsection{Exact cost of the read prefix with arbitrary contents}

Attach any graph $H$ to the local source's content output. Transport the
read head through serial composition and compare an arbitrary actual step
with the serially embedded original local read. Dereference coherence shows
that every successor is isomorphic to the explicit result retaining $H$.
Neither termination, stability nor acyclicity of $H$ is assumed.

\lean{Copying/DereferenceCost}{9}{15}{Dereference.contents_successor}

\termdef{read-cost-handoff}
Let $s$ be this composed graph with a search token at its request, and let
$p$ be the composed original reduct with a rewrite token at that request.
The \emph{read cost handoff} is
\[
 s\text{ terminates in }n
 \quad\Longleftrightarrow\quad
 \exists r,\ n=r+5\ \land\ p\text{ terminates in }r.
\]
Its five transitions are: search the instance argument; return from the
instance; request dereference; perform the store-preserving read rewrite;
request copying at the new distributor. Apply the exact search-cost theorem
at the four search transitions and the dereference cost theorem at the
computation. Eliminating the five intermediate natural numbers gives the
displayed equation; conversely, inserting $r+4,r+3,r+2,r+1,r$ constructs
all five predecessor costs. This is a full-machine equivalence for all
successful costs, not an existence-only local execution.

\lean{Copying/DereferenceCost}{18}{38}{Dereference.contents_cost_handoff}

\Needspace{12\baselineskip}
The five-step prefix itself always exists: embed the previously proved
local read run into the combined machine, then transport it through serial
attachment. No successful execution can have cost at most five. If it did,
the handoff equation would force the residual cost to be zero. A zero-step
successful run would make $p$ final, but its mode is rewrite, whereas finality
requires backtrack at an input. This proves that the mandatory pending copy
has not been erased, even for literal data.

\lean{Copying/DereferenceCost}{41}{45}{Dereference.contents_prefix}

\lean{Copying/DereferenceCost}{48}{56}{Dereference.contents_no_early_success}

Successor coherence also holds after inserting the read into any typed
host hole, transporting interfaces along their equalities. Transport the
head through rebounding and plugging, then compare an arbitrary step with
the explicitly plugged original local read. This permits internal resumed
request wires. The fixed five-step cost statement above concerns the serial
content attachment; an arbitrary surrounding host can perform further work.

\lean{Copying/DereferenceCost}{59}{68}{Dereference.host_successor}

\Needspace{12\baselineskip}
\subsection{Unchanged translated data and depth padding}

The closed program \texttt{storedTerm} places the literal Figure~16
translation of any closed term $t$ at the read source's content output.
Its depth is the translation depth of $t$; focus and hole freedom follow
from the safety of the read source, translation safety and serial composition.
The state \texttt{pending} is the explicit original read reduct with the
same translation attached and its token in rewrite mode. These are heap-state
regressions; no identification with a whole surface-language read translation
is asserted.

\lean{Copying/DereferenceCoherenceExamples}{9}{14}{DereferenceCoherenceExamples.storedTerm}

\lean{Copying/DereferenceCoherenceExamples}{17}{21}{DereferenceCoherenceExamples.pending}

For every finite depth padding and every proposed successful cost, the
unbounded initial execution satisfies exactly the read cost handoff to
\texttt{pending}. First use the proved padding-invariance theorem to reflect
successful costs to the original allowance. Then use the bounded/unbounded
successful-run correspondence and instantiate the arbitrary-contents handoff.
Lifting the serial prefix also proves an actual five-step unbounded run to
the specified pending state. The same zero-residual-cost contradiction proves
that no padding permits a successful cost at most five.

\lean{Copying/DereferenceCoherenceExamples}{24}{31}{DereferenceCoherenceExamples.storedTerm_cost_handoff}

\lean{Copying/DereferenceCoherenceExamples}{34}{39}{DereferenceCoherenceExamples.storedTerm_prefix}

\lean{Copying/DereferenceCoherenceExamples}{42}{50}{DereferenceCoherenceExamples.storedTerm_no_early_success}

Specializing the general cost law to a natural-number literal retains its
complete translation and copying continuation. Specializing it to the literal
self-application term $\Omega$ retains its nested boxes as stored data;
this proves no separate divergence or termination result for $\Omega$.
Finally, arbitrary-successor coherence specializes to every translated payload,
preserving the entire attached graph and resumed token under every actual
contextual read.

\lean{Copying/DereferenceCoherenceExamples}{53}{59}{DereferenceCoherenceExamples.literal_cost_handoff}

\lean{Copying/DereferenceCoherenceExamples}{62}{67}{DereferenceCoherenceExamples.omega_cost_handoff}

\lean{Copying/DereferenceCoherenceExamples}{70}{77}{DereferenceCoherenceExamples.storedTerm_successor}
